\documentclass[runningheads]{llncs}

\usepackage[T1]{fontenc}
\usepackage{graphicx}
\usepackage{amsmath}
\usepackage{amssymb}
\usepackage{booktabs}
\usepackage{multirow}
\usepackage{xcolor}
\usepackage{placeins}
\usepackage[hidelinks]{hyperref}
\renewcommand\UrlFont{\color{blue}\rmfamily}
\urlstyle{rm}

% Placeholder box so the document always compiles before final art lands.
\newcommand{\figph}[2]{%
  \fbox{\begin{minipage}[c][#1][c]{\dimexpr\linewidth-2\fboxsep-2\fboxrule}
    \centering\footnotesize\color{gray}#2
  \end{minipage}}}

% Notation.
\newcommand{\LGE}{\mathrm{LGE}}
\newcommand{\Cz}{\mathrm{C0}}
\newcommand{\TT}{\mathrm{T2}}
\newcommand{\Img}{\mathbf{I}}
\newcommand{\sg}{\mathrm{sg}}
\newcommand{\warp}{\mathcal{W}}

% numeric macros formerly generated by scripts/make_paper_tables.py
\newcommand{\edemaZoneFull}{0.634\,{\tiny $\pm$\,0.013}}
\newcommand{\edemaZoneNoCz}{0.632\,{\tiny $\pm$\,0.013}}
\newcommand{\edemaZoneDirect}{0.349\,{\tiny $\pm$\,0.070}}

\begin{document}

\title{MoSAIC: Motion- and Supervision-Aware Inference for\\
Myocardial Scar and Edema Segmentation from\\
Multi-Sequence and Cine CMR}
\titlerunning{MoSAIC}

\author{Anonymous Author(s)}
\authorrunning{Anonymous}
\institute{Institution masked for double-blind review}

\maketitle

\begin{abstract}
Segmenting myocardial scar and edema from cardiac magnetic resonance (CMR)
quantifies infarct burden and salvageable myocardium, while realistic
multi-center data require models that use heterogeneous evidence. Sequence
availability and label spaces vary across centers, so supervision should follow
the annotated class set of each case. Cine CMR adds complementary motion
evidence: scar appears mainly as regional wall-motion abnormality. We propose
MoSAIC (Motion- and Supervision-Aware Inference for Myocardial Scar and Edema
Segmentation), a coarse-to-fine, heterogeneity-aware, anatomy-conditioned cascade
for multi-sequence and cine CMR. A coarse stage first estimates cardiac anatomy
and defines a spatial prior. The fine stage then combines a multi-sequence
scar/anatomy decoder, a dedicated edema-zone expert and a cine motion scar head.
Presence-gated encoding with a supervision-masked objective ensures that each
case supervises only its annotated classes, allowing heterogeneous cohorts to be
pooled with case-aware supervision. The cine pathology head uses a dense
end-diastole-anchored motion field as kinematic evidence, and the edema expert
predicts the injury-zone label used for scar/edema fusion. Official validation on
CARE-2026 CARE-Myocardium shows best MoSAIC Dice entries of 0.6965 for MyoPS
scar, 0.5983 for MyoPS edema and 0.2069 for CineMyoPS scar. Internal five-fold
evaluation reports Dice of 0.668 for scar and 0.645 for lesion union on
multi-sequence CMR, and 0.457 for cine-only scar.

\keywords{Myocardial pathology segmentation \and Multi-sequence CMR \and
Cine CMR \and Motion-derived evidence \and Partial supervision.}
\end{abstract}

% ===========================================================================
\section{Introduction}
% ===========================================================================

After acute myocardial infarction, scar marks irreversibly necrotic myocardium
while edema marks salvageable injured tissue. Their relation
guides reperfusion assessment~\cite{eitel2011cardiac,ibanez2015evolving}, and
CMR resolves them through complementary contrasts: LGE highlights
scar~\cite{kim1999relationship}, T2-weighted imaging highlights
edema~\cite{abdelaty2004delayed}, and bSSFP/C0 supplies anatomy. Automating
their joint delineation is the goal of the MyoPS benchmark
family~\cite{li2023myops,zhuang2019multivariate}.

The CARE-2026 setting emphasizes realistic heterogeneity in both image evidence
and annotation protocols. In multi-sequence CMR, training cases may contain LGE,
LGE+C0 or the full LGE+C0+T2 combination; 80 of 220 cases annotate edema, and
centers provide different anatomical label sets. Effective training therefore
requires the objective to respect the class set that each center annotated. In
cine-only CMR, scar is represented primarily through regional wall-motion
abnormality~\cite{zhang2019deeplearningcine,ding2025cinemyops}, making motion a
clinically meaningful signal for contrast-free pathology inference.

We propose a coarse to fine heterogeneity-aware framework that models both
motion evidence and supervision availability (Fig.~\ref{fig:pipeline}).
The coarse stage predicts cardiac anatomy, extracts a region of interest and
builds a spatial prior. Inside this prior, Stage~2a decodes multi-sequence scar
with auxiliary anatomy consistency, Stage~2b predicts the challenge injury zone
(edema $\cup$ scar), and Stage~2c derives cine scar from ED-anchored motion and
warped anatomy. A supervision-masked loss ensures that each case supervises
only its annotated classes, enabling heterogeneous centers to be pooled while
preserving their annotation semantics.

Our contributions are a coarse to fine, heterogeneity-aware cascade that links
anatomy localization with scar/anatomy and edema-zone decoding; a
presence-gated, supervision-masked recipe for ragged multi-center labels; a
motion-driven cine scar head; and a five-fold CARE-2026 evaluation with
organizer-scored validation results.

% ===========================================================================
\section{Related Work}
% ===========================================================================

The MyoPS challenges established joint scar/edema segmentation on paired LGE,
T2 and bSSFP data~\cite{zhuang2019multivariate,li2023myops}. Subsequent methods
improved flexible sequence use~\cite{qiu2023myopsnet} and joint
alignment-segmentation~\cite{ding2023aligning}, with a primary focus on unified
annotation spaces. Generic U-Net and nnU-Net pipelines provide strong
segmentation backbones~\cite{ronneberger2015unet,isensee2021nnunet}; MyoPS
variants have also used coarse cardiac localization before fine scar/edema
decoding~\cite{zhai2020coarsemyops,wang2022awsnet}. The proposed framework
follows this anatomy-prior intuition and extends it to variable labels, variable
sequence availability and cine motion evidence.

Missing-modality methods such as HeMIS handle variable inputs through pooled or
imputed latent statistics~\cite{havaei2016hemis,dorent2019hetero}, while
partial-label losses preserve the semantics of classes outside a case's
annotation set~\cite{shi2021marginal,zhou2019priorpartial,fidon2021labelsetloss}.
The proposed framework combines the two requirements in one mechanism:
modality presence gates and a per-class supervision mask in the loss normalizer.

Contrast-free cine CMR offers a complementary route to pathology inference.
Prior work has used cine appearance~\cite{zhang2019deeplearningcine},
strain~\cite{scatteia2017strain,%
morales2021deepstrain}, or joint motion-and-segmentation
learning~\cite{qin2018jointmotion}; CineMyoPS aggregates motion and anatomy
branches across the cardiac cycle~\cite{ding2025cinemyops}. The proposed
framework follows the same clinical motivation and feeds the scar head
motion-derived evidence, emphasizing the kinematic signal most directly tied to
cine CMR.

% ===========================================================================
\section{Preliminaries}
% ===========================================================================

\noindent\textbf{Problem setup.} A dataset is a set of tuples
$\mathcal{D}=\{(X_i, Y_i, \mathbf{m}_i, \mathbf{s}_i)\}$. The label map
$Y \in \{0,\dots,C\}^{D\times H\times W}$ uses the challenge classes
$\{$background, myocardium, LV, RV, edema, scar$\}$. For label fusion, we denote
the injury zone as $Z=\text{edema}\cup\text{scar}$ and recover pure edema after
scar prediction. Two acquisition regimes occur. In the
\emph{multi-sequence} regime $X=(\Img_\LGE,\Img_\Cz,\Img_\TT)$ and
$\mathbf{m}\in\{0,1\}^3$ records which sequences a center actually acquired,
with absent sequences zero-filled. In the \emph{cine} regime
$X=(I_0,\dots,I_{T-1})$ is a single bSSFP sequence sampled over the cardiac
cycle and anchored at end-diastole (ED). In both regimes
$\mathbf{s}\in\{0,1\}^{C}$ records which classes the originating center
annotated; $s_c=0$ means the label map is silent about class $c$, leaving its
underlying presence unspecified.

\noindent\textbf{Evaluation.} Internal cross-validation reports the Dice
similarity coefficient and the 95th-percentile Hausdorff distance (HD95, in mm)
computed in the resampled physical space. Organizer-scored validation uses the
CARE-Myocardium metrics, Dice and Hausdorff distance (HD, in mm).

% ===========================================================================
\section{Method}
% ===========================================================================

\begin{figure}[t]
\centering
\includegraphics[width=\linewidth]{figures/fig1_original_redraw.pdf}
\caption{Overview of MoSAIC. A coarse anatomy stage defines the spatial prior;
separate multi-sequence scar, edema-zone, and cine-motion branches produce
task-specific predictions that are fused in the challenge label space.}
\label{fig:pipeline}
\end{figure}

The framework follows a coarse to fine design. Stage~1 performs coarse cardiac
anatomy localization and produces a spatial prior. Stage~2 splits fine decoding
into a multi-sequence scar/anatomy expert, a dedicated edema-zone expert and a
cine motion scar expert. We describe each in turn.

\subsection{Stage 1: Anatomy Conditioning}

A compact 2D U-Net~\cite{ronneberger2015unet} $F_{\text{coarse}}$ predicts three
coarse classes (myocardial ring, LV, RV) from the channel-stacked sequences
augmented with constant \emph{presence maps}, i.e.\ the input is
$[\Img_\LGE,\Img_\Cz,\Img_\TT, m_1\mathbf{1}, m_2\mathbf{1}, m_3\mathbf{1}]$.
The presence maps let one network serve every acquisition protocol: the decoder
can condition on which channels are informative and avoids interpreting a
zero-filled channel as dark tissue. Each output channel is thresholded and
reduced to its largest connected component, giving
$\hat{Y}^{\text{c}}$ and the binary prior $P=\mathbf{1}\!\left[\hat{Y}^{\mathrm{c}}>0\right]$.
The bounding box of $P$, dilated by a fixed margin, defines the region of
interest $\Omega$ passed to Stage~2. Because $P$ is derived from anatomy rather
than pathology, it can be computed for cases with or without lesion annotations
and is used as the common conditioning signal for the fine stage.

\subsection{Stage 2a: Multi-Sequence Scar and Anatomy Decoding}

\noindent\textbf{Role.} Stage~2a is the multi-sequence scar/anatomy expert. It
predicts scar from LGE-centered fused features and uses auxiliary myocardium
decoding to keep the fused representation anatomy-aware; edema-zone decoding is
handled by Stage~2b.

\noindent\textbf{Modality-decoupled encoding with presence gating.} The framework
runs three independent encoders
$E^{(m)}$, $m\in\{\LGE,\Cz,\TT\}$, each producing a four-level pyramid
$f^{(m)}_\ell$, $\ell=1,\dots,4$. Missing sequences are suppressed
multiplicatively,
\begin{equation}
\tilde{f}^{(m)}_\ell \;=\; m_m \cdot f^{(m)}_\ell ,
\end{equation}
so an absent auxiliary sequence contributes no image-derived feature maps; the
decoder still receives fixed sigmoid gate values from the zeroed features. The
same weights serve LGE-only, LGE+C0 and full three-sequence cases.

\noindent\textbf{An optional alignment branch.} Separate breath-holds can
introduce fine non-rigid offsets around the thin myocardial ring where lesion
decisions are made. The framework therefore includes a deformable branch that models
sequence alignment inside the representation the decoder consumes: for each moving
sequence $m\in\{\Cz,\TT\}$ a localization head $h^{(m)}$ reads the \emph{paired}
bottlenecks and predicts $K=g^2$ control-point offsets
$\theta^{(m)}=h^{(m)}([\tilde{f}^{(m)}_4, f^{(\LGE)}_4])$, defining a thin-plate
spline~\cite{bookstein1989principal} $\mathcal{T}_{\theta}$ that warps every
pyramid level of that sequence (Appendix~\ref{app:tps}). The head is
zero-initialized for stable alignment learning under the segmentation objective.

\noindent\textbf{Fusion and prior gating.} The decoder fuses the sequence
pyramids with a gating rule that keeps LGE as the structural carrier while the
other sequences modulate it, and the Stage-1 prior re-enters at every decoder
level as a multiplicative gate that emphasizes the cardiac region while retaining
context. Appendix~\ref{app:fusion} gives the equations and the fixed gate
values induced by absent auxiliary sequences.

\subsection{A Supervision-Masked Objective for Ragged Annotation}

Training pools cases with different annotated class subsets by applying the
per-class supervision mask $\mathbf{s}$ directly in the loss normalizer. With $\ell^{\text{dice}}_{b,c}$ and
$\ell^{\text{bce}}_{b,c}$ the per-sample per-class Dice and binary
cross-entropy terms and $w_c$ a class weight,
\begin{equation}
\mathcal{L}_{\text{seg}}
=\frac{\sum_{b,c}\,s_{b,c}\,w_c\left[\ell^{\text{dice}}_{b,c}+\ell^{\text{bce}}_{b,c}\right]}
       {\sum_{b,c}\,s_{b,c}} .
\label{eq:masked}
\end{equation}
Each annotated class contributes its supervised gradient; an unannotated class
receives zero weight in both the numerator and the denominator. A case from a
scar-focused center therefore supervises the scar head and leaves edema
supervision to centers that annotated edema.

Small lesions additionally receive a recall-oriented Tversky
term~\cite{salehi2017tversky} with $\alpha<\beta$ and a focal
term~\cite{lin2017focal}, both restricted to lesion channels and both masked the
same way. Deep supervision is applied at two decoder levels.

An auxiliary \emph{cross-sequence myocardium consistency} term completes
Stage~2a (Appendix~\ref{app:losses}). It decodes a myocardium mask from each
sequence-specific feature stack and regularizes all available pairs with Dice
agreement. This ties representations through shared anatomy without adding
annotations.

\subsection{Stage 2b: A Dedicated Edema Expert}

Scar and edema are modeled with dedicated experts because they rest on different
physical contrasts, LGE hyper-enhancement versus T2 hyper-intensity, and receive
different supervision density: in this cohort scar is annotated on 220 cases and
edema on 80. A separate edema expert lets the T2-driven target use a focused
decoder and objective.

The framework therefore trains a separate expert. T2 forms the main stream while the
LGE and C0 streams are fused element-wise into a \emph{reference} pyramid
$e^{\text{ref}}_\ell$, which enters the decoder \emph{detached},
$d_\ell=\Psi_\ell([\mathrm{up}(d_{\ell+1}),\,\sg(e^{\text{ref}}_\ell),\,e^{\TT}_\ell])$.
Anatomical context therefore conditions the decision without letting the edema
gradient reshape the LGE/C0 features that the scar expert relies on.

The expert regresses the injury zone $Z=\text{edema}\cup\text{scar}$, matching
the label representation used for fusion. Pure edema is recovered at fusion time
by assigning the predicted scar first and then filling
$Z\setminus\widehat{\text{scar}}$ as edema.

\subsection{Stage 2c: Motion-Derived Scar from Cine}

For cine-only data, the framework replaces the multi-sequence encoder stack with a
motion extractor. A shared encoder embeds each of $T$ ED-anchored frames; an
anatomy decoder predicts per-frame myocardium and LV; and a motion decoder
predicts a dense displacement field between the ED reference $I_0$ and frame
$I_t$, bounded to keep it physically plausible:
\begin{equation}
\varphi_t=\tau\tanh\!\left(D_{\text{mot}}(f_0,f_t)\right),\qquad \varphi_0=\mathbf{0},
\end{equation}
with $\tau$ the maximum displacement as a fraction of the field of view. The
field warps both the anatomy probability and the frame,
$\hat{A}_t=\warp(\sigma(A_t),\varphi_t)$ and $\hat{I}_t=\warp(I_t,\varphi_t)$.

The pathology head receives only motion, warped anatomy and the anatomical prior:
\begin{equation}
z_t=\Theta\!\left(\left[\varphi_t,\;\hat{A}_t,\;P\right]\right),
\qquad
\hat{Y}=\mathrm{Conv}_{1\times1}\!\left(\tfrac{1}{T}\textstyle\sum_{t} z_t\right).
\label{eq:cinehead}
\end{equation}
This restricts the cine scar branch to contrast-free temporal and anatomical
cues.

Training combines the masked pathology loss with four auxiliary terms: ED-frame
anatomy supervision; consistency between each warped anatomy map and the ED
prediction; a photometric term matching $\hat{I}_t$ to $I_0$; and an $\ell_1$
smoothness penalty on $\nabla\varphi$ for the displacement field.

\subsection{Inference}

Both regimes share one skeleton: predict anatomy, restrict to the prior, predict
pathology with flip test-time augmentation, then apply anatomical containment and
minimum lesion-size cleanup.
For the multi-sequence regime, scar comes from the Stage-2a scar/anatomy expert,
and edema is assigned from $Z\setminus\widehat{\text{scar}}$ under the injury-zone
fusion convention. For cine, predictions are additionally averaged over three
through-plane resamplings as the deployment-time handling of through-plane
resolution variation.

% ===========================================================================
\section{Experiments}
\label{sec:exp}
% ===========================================================================

\noindent\textbf{Data and data-use declaration.} All experiments use the
CARE-2026 CARE-Myocardium dataset, comprising the MyoPS multi-sequence cohort
and the CineMyoPS cine cohort. We use no other data. We gratefully acknowledge
the challenge organizers, and cite the sources as required:
\cite{zhuang2019multivariate,qiu2023myopsnet,ding2023aligning,ding2025cinemyops}.
Table~\ref{tab:cohort} summarises the heterogeneity that motivates our design.
No pre-trained weights are used; every model is trained from scratch.

\begin{table}[t]
\centering
\caption{The CARE-2026 CARE-Myocardium training cohort. Sequence availability and
annotated classes both vary by center: 116 of 220 multi-sequence cases provide LGE
alone, and 80 carry edema labels. The evaluation cohort provides all three
sequences, so a model can learn from partial training protocols and exploit
complete inputs at test time.}
\label{tab:cohort}
\setlength{\tabcolsep}{5pt}
\begin{tabular}{llll}
\toprule
Centre & Cases & Sequences & Annotated classes \\
\midrule
\multicolumn{4}{l}{\textit{MyoPS (multi-sequence)}}\\
A & 81 & LGE & myo, LV, scar \\
B & 35 & LGE+C0+T2 & myo, LV, RV, edema, scar \\
C & 45 & LGE+C0+T2 & myo, LV, RV, edema, scar \\
E & 7 & LGE+C0 & myo, LV, RV, scar \\
F & 9 & LGE+C0 & myo, LV, RV, scar \\
G & 8 & LGE+C0 & myo, LV, RV, scar \\
H & 35 & LGE & scar \\
\midrule
\multicolumn{4}{l}{\textit{CineMyoPS (cine only)}}\\
$\alpha$ & 40 & Cine & myo, LV, scar \\
$\beta$ & 24 & Cine & myo, LV, scar \\
\bottomrule
\end{tabular}
\end{table}

\noindent\textbf{Implementation.} Volumes are resampled to
$1.25\times1.25\times10.0$\,mm (multi-sequence) and
$1.25\times1.25\times8.0$\,mm (cine), robustly z-scored per volume, and processed
slice-wise at $192\times192$. Stage-1 uses 24 base channels, Stage-2a 24, the
edema expert 64 and the cine model 16, with instance
normalization~\cite{ulyanov2016instance}. We train with
AdamW~\cite{loshchilov2019adamw} under a cosine schedule~\cite{loshchilov2017sgdr},
gradient clipping at 0.1, and a class-quota batch sampler that guarantees a fixed
fraction of lesion-bearing slices per batch. The TPS warm-up is $E_w=10$ epochs.
Cross-validation uses five folds stratified by center and by lesion presence, so
every fold preserves the cohort's modality and annotation mix. Because edema is
annotated at two centers, edema metrics are computed on those centers'
cases. All runs use seed 3407.

\noindent\textbf{Metrics.} Internal five-fold results are reported as Dice
and HD95 (mm), with mean\,$\pm$\,standard deviation across folds. Official
validation results are organizer-reported Dice and HD (mm).

\subsection{Main Results}

\begin{table}[t]
\centering
\caption{Five-fold cross-validated framework performance, mean\,$\pm$\,s.d.\
across folds. Pure-edema metrics use cases with edema annotations; the remaining
metrics use their corresponding annotated cohorts. Cine scar is inferred from
motion alone, without any contrast agent.}
\label{tab:main}
\begin{tabular}{lcc}
\toprule
Structure & Dice $\uparrow$ & HD95 (mm) $\downarrow$ \\
\midrule
\multicolumn{3}{l}{\textit{MyoPS: multi-sequence (LGE, C0, T2)}}\\
\quad Myocardium & 0.687\,{\tiny $\pm$\,0.036} & 8.29\,{\tiny $\pm$\,1.61} \\
\quad LV & 0.866\,{\tiny $\pm$\,0.015} & 5.87\,{\tiny $\pm$\,1.43} \\
\quad RV & 0.841\,{\tiny $\pm$\,0.022} & 6.63\,{\tiny $\pm$\,1.80} \\
\quad Pure edema & 0.349\,{\tiny $\pm$\,0.065} & 20.21\,{\tiny $\pm$\,3.49} \\
\quad Scar & 0.668\,{\tiny $\pm$\,0.022} & 11.69\,{\tiny $\pm$\,1.06} \\
\quad Lesion union & 0.645\,{\tiny $\pm$\,0.024} & 15.87\,{\tiny $\pm$\,3.13} \\
\midrule
\multicolumn{3}{l}{\textit{CineMyoPS: cine only}}\\
\quad Myocardium & 0.695\,{\tiny $\pm$\,0.024} & 5.35\,{\tiny $\pm$\,0.30} \\
\quad LV & 0.919\,{\tiny $\pm$\,0.007} & 3.87\,{\tiny $\pm$\,0.55} \\
\quad Scar & 0.457\,{\tiny $\pm$\,0.038} & 19.68\,{\tiny $\pm$\,3.88} \\
\bottomrule
\end{tabular}
\end{table}

\noindent\textbf{Multi-sequence segmentation.} The upper block of
Table~\ref{tab:main} reports five-fold performance on the MyoPS cohort.
Anatomy shows high Dice: LV reaches $0.866$ and RV $0.841$, with
myocardium at $0.687$ for the thin ring-shaped structure. Among pathology
classes, scar reaches $0.668$ Dice and $11.69$\,mm HD95. The scar branch uses
LGE as its primary contrast. The lesion union reaches $0.645$ Dice and is
reported in the same injury-zone label space used for fusion; pure edema is
reported as a local two-center T2 endpoint in
Table~\ref{tab:main}.

\noindent\textbf{Cine-only segmentation.} The cine block shows strong anatomy
(LV $0.919$, myocardium $0.695$) on the two-center CineMyoPS cohort. The
motion-based cine branch reaches $0.457$ Dice with an HD95 of $19.68$\,mm using
contrast-free temporal input.

\noindent\textbf{Official validation leaderboard.} Table~\ref{tab:leaderboard}
reports the best MoSAIC organizer-scored CARE-2026 validation entries for each
official target, using the organizer metrics Dice and HD. The best MoSAIC entries
reach $0.6965$ Dice for MyoPS scar, $0.5983$ for MyoPS edema and $0.2069$ for
CineMyoPS scar. Compared with the nnU-Net comparator, MoSAIC improves MyoPS scar,
CineMyoPS scar, and the mean Dice over the three official targets ($0.492$ to
$0.501$).

\begin{table}[t]
\centering
\caption{Best organizer-scored MoSAIC entries and the nnU-Net comparator on the
CARE-2026 validation leaderboards, reported as Dice/HD (mm). For each target, the
MoSAIC row is the submitted entry with the highest Dice.}
\label{tab:leaderboard}
\scriptsize
\begin{tabular}{llcccc}
\toprule
Leaderboard & Target & MoSAIC Dice $\uparrow$ & MoSAIC HD $\downarrow$ & nnU-Net Dice $\uparrow$ & nnU-Net HD $\downarrow$ \\
\midrule
MyoPS & scar & 0.6965 & 13.7827 & 0.6258 & 14.9844 \\
MyoPS & edema & 0.5983 & 26.7067 & 0.6691 & 21.0898 \\
CineMyoPS & scar & 0.2069 & 48.7463 & 0.1816 & 52.9706 \\
\bottomrule
\end{tabular}
\end{table}

% Temporarily omitted while tightening the paper.
\iffalse
\subsection{Ablation Study}

\begin{table}[t]
\centering
\caption{MyoPS ablation, five-fold cross-validated Dice. Every variant is
retrained from scratch on the same folds and consumes the same frozen Stage-1
prior, so Stage-1 is fixed across all comparisons. $\Delta$ reports scar Dice
relative to the full model.}
\label{tab:ablation_myops}
\setlength{\tabcolsep}{4pt}
\begin{tabular}{lcccc}
\toprule
Variant & Scar & Edema & Lesion union & $\Delta$Scar \\
\midrule
\textbf{Full framework} & 0.668\,{\tiny $\pm$\,0.022} & 0.349\,{\tiny $\pm$\,0.065} & 0.645\,{\tiny $\pm$\,0.024} & -- \\
\midrule
\multicolumn{5}{l}{\textit{Architecture}} \\
\quad w/o TPS feature alignment & 0.663\,{\tiny $\pm$\,0.023} & 0.345\,{\tiny $\pm$\,0.067} & 0.646\,{\tiny $\pm$\,0.022} & \textcolor{gray}{-0.005} \\
\quad early fusion (single shared encoder) & 0.652\,{\tiny $\pm$\,0.026} & 0.332\,{\tiny $\pm$\,0.057} & 0.645\,{\tiny $\pm$\,0.021} & \textcolor{red}{-0.016} \\
\quad w/o Stage-1 anatomy cascade & 0.638\,{\tiny $\pm$\,0.029} & 0.354\,{\tiny $\pm$\,0.066} & 0.644\,{\tiny $\pm$\,0.023} & \textcolor{red}{-0.030} \\
\quad w/o spatial prior gate & 0.636\,{\tiny $\pm$\,0.024} & 0.356\,{\tiny $\pm$\,0.070} & 0.644\,{\tiny $\pm$\,0.024} & \textcolor{red}{-0.032} \\
\quad w/o dedicated edema expert & 0.668\,{\tiny $\pm$\,0.022} & 0.137\,{\tiny $\pm$\,0.030} & 0.275\,{\tiny $\pm$\,0.033} & \textcolor{gray}{+0.000} \\
\midrule
\multicolumn{5}{l}{\textit{Objective}} \\
\quad w/o supervision masking & 0.658\,{\tiny $\pm$\,0.015} & 0.343\,{\tiny $\pm$\,0.069} & 0.642\,{\tiny $\pm$\,0.024} & \textcolor{red}{-0.010} \\
\quad w/o cross-sequence consistency & 0.650\,{\tiny $\pm$\,0.028} & 0.349\,{\tiny $\pm$\,0.070} & 0.644\,{\tiny $\pm$\,0.023} & \textcolor{red}{-0.018} \\
\quad w/o lesion Tversky\,/\,focal terms & 0.669\,{\tiny $\pm$\,0.019} & 0.347\,{\tiny $\pm$\,0.062} & 0.646\,{\tiny $\pm$\,0.022} & \textcolor{gray}{+0.001} \\
\midrule
\multicolumn{5}{l}{\textit{Inference}} \\
\quad w/o test-time augmentation & 0.659\,{\tiny $\pm$\,0.019} & 0.345\,{\tiny $\pm$\,0.062} & 0.643\,{\tiny $\pm$\,0.020} & \textcolor{red}{-0.009} \\
\quad w/o anatomical constraint & 0.668\,{\tiny $\pm$\,0.021} & 0.346\,{\tiny $\pm$\,0.064} & 0.642\,{\tiny $\pm$\,0.023} & \textcolor{gray}{+0.000} \\
\quad w/o component cleanup & 0.664\,{\tiny $\pm$\,0.026} & 0.348\,{\tiny $\pm$\,0.065} & 0.645\,{\tiny $\pm$\,0.022} & \textcolor{gray}{-0.004} \\
\bottomrule
\end{tabular}
\end{table}

\begin{table}[t]
\centering
\caption{CineMyoPS ablation, five-fold cross-validated Dice. Temporal aggregation
is the largest factor in this cohort; motion remains a useful inductive bias for
cine-only scar.}
\label{tab:ablation_cine}
\setlength{\tabcolsep}{4pt}
\begin{tabular}{lcccc}
\toprule
Variant & Scar & Myocardium & LV & $\Delta$Scar \\
\midrule
\textbf{Full framework} & 0.457\,{\tiny $\pm$\,0.038} & 0.695\,{\tiny $\pm$\,0.024} & 0.919\,{\tiny $\pm$\,0.007} & -- \\
\midrule
\multicolumn{5}{l}{\textit{Evidence source}} \\
\quad raw intensity evidence & 0.464\,{\tiny $\pm$\,0.034} & 0.703\,{\tiny $\pm$\,0.021} & 0.919\,{\tiny $\pm$\,0.007} & \textcolor{teal}{+0.006} \\
\quad single ED frame ($T{=}1$) & 0.420\,{\tiny $\pm$\,0.035} & 0.697\,{\tiny $\pm$\,0.018} & 0.920\,{\tiny $\pm$\,0.007} & \textcolor{red}{-0.037} \\
\midrule
\multicolumn{5}{l}{\textit{Auxiliary objectives}} \\
\quad w/o flow smoothness & 0.456\,{\tiny $\pm$\,0.029} & 0.712\,{\tiny $\pm$\,0.009} & 0.920\,{\tiny $\pm$\,0.007} & \textcolor{gray}{-0.002} \\
\quad w/o temporal anatomy consistency & 0.459\,{\tiny $\pm$\,0.044} & 0.693\,{\tiny $\pm$\,0.033} & 0.919\,{\tiny $\pm$\,0.007} & \textcolor{gray}{+0.001} \\
\midrule
\multicolumn{5}{l}{\textit{Inference}} \\
\quad w/o test-time augmentation & 0.452\,{\tiny $\pm$\,0.037} & 0.687\,{\tiny $\pm$\,0.025} & 0.918\,{\tiny $\pm$\,0.007} & \textcolor{red}{-0.006} \\
\quad w/o anatomical constraint & 0.457\,{\tiny $\pm$\,0.038} & 0.695\,{\tiny $\pm$\,0.024} & 0.919\,{\tiny $\pm$\,0.007} & \textcolor{gray}{+0.000} \\
\quad w/o component cleanup & 0.451\,{\tiny $\pm$\,0.035} & 0.695\,{\tiny $\pm$\,0.024} & 0.919\,{\tiny $\pm$\,0.007} & \textcolor{red}{-0.006} \\
\bottomrule
\end{tabular}
\end{table}

\noindent\textbf{Edema expert zone formulation.} Trained as a
zone regressor with the LGE+C0 reference stream, the expert reaches
\edemaZoneFull\ Dice on the injured territory. The LGE-only reference variant
gives \edemaZoneNoCz, and direct annulus regression gives
\edemaZoneDirect. The comparison isolates which part of the design carries the
performance. The C0 reference stream and LGE-only reference variant remain
within 0.002 Dice, one standard deviation across folds. The zone formulation is
decisive: predicting the simply connected injured territory reaches
\edemaZoneFull\ Dice with a compact fold-to-fold spread of $\pm$0.013, while
the direct annulus formulation gives \edemaZoneDirect. The zone target therefore
provides the stable supervision used by the edema expert.

\noindent\textbf{Motion evidence for cine.}
Table~\ref{tab:ablation_cine} reveals a more nuanced picture than the
architectural stance of Eq.~\eqref{eq:cinehead} might suggest. Replacing the
motion field with raw frame intensity at the pathology head changes scar Dice by
$+0.006$, well within the fold-to-fold spread of $\pm 0.038$: on this cohort,
the two representations carry comparable information about regional dysfunction.
\emph{Temporal aggregation} provides the clearest signal: the full cardiac-cycle
model reaches $0.457$ Dice, compared with $0.420$ for the single-ED-frame
variant. Together these results suggest that temporal aggregation is useful on
this cohort; whether the motion representation improves scanner transfer remains
a separate validation question.

\noindent\textbf{Does presence gating support sequence variation?}
Table~\ref{tab:modality} evaluates the \emph{same} trained model under
controlled sequence variants at inference time. The full LGE+C0+T2 protocol
achieves the strongest balanced anatomy and pathology scores. The LGE+C0 variant
keeps scar at $0.620$ Dice, and the LGE-only variant keeps scar at $0.613$ Dice,
showing that LGE carries the dominant pathological signal while C0 contributes
anatomical detail. These comparisons confirm that the multiplicative presence
gate adapts gracefully across acquisition protocols.

\begin{table}[t]
\centering
\caption{Robustness to sequence availability. The \emph{same} trained model is
evaluated under controlled acquisition protocols by clearing the corresponding
presence flags and zero-filling their channels.}
\label{tab:modality}
\begin{tabular}{lccc}
\toprule
Sequences available & Scar & Edema & Myocardium \\
\midrule
LGE + C0 + T2 (full) & 0.668\,{\tiny $\pm$\,0.022} & 0.349\,{\tiny $\pm$\,0.065} & 0.687\,{\tiny $\pm$\,0.036} \\
LGE + C0 & 0.620\,{\tiny $\pm$\,0.028} & 0.197\,{\tiny $\pm$\,0.042} & 0.606\,{\tiny $\pm$\,0.055} \\
LGE only & 0.613\,{\tiny $\pm$\,0.032} & 0.162\,{\tiny $\pm$\,0.044} & 0.523\,{\tiny $\pm$\,0.102} \\
\bottomrule
\end{tabular}
\end{table}

\begin{figure}[t]
\centering
\includegraphics[width=\textwidth]{figures/fig2_qualitative}
\caption{Qualitative multi-sequence results on the two best-scoring cases, each
predicted by the model of the fold that held it out. Rows are cases, columns walk
from the three input sequences to the ground truth and the framework output.}
\label{fig:qual}
\end{figure}
\fi

\FloatBarrier
\subsection{Discussion}

\noindent\textbf{Summary.} Coarse anatomy and case-aware supervision provide a
common interface for heterogeneous centers. Stage~1 supplies a shared spatial
frame, Stage~2a and Stage~2b separate scar and edema according to their evidence
sources, and Stage~2c substitutes temporal motion for contrast in cine CMR. In
official validation, the strongest MoSAIC entries improve MyoPS scar and
CineMyoPS scar relative to the nnU-Net comparator, while the three-target mean
Dice is slightly higher.

\noindent\textbf{Limitations.} Cine folds are drawn from the same source centers,
whereas official validation may change scanner, frame count, temporal resolution
and slice thickness; protocol-normalized flow or explicit frame-rate augmentation
is therefore a natural extension. For edema, the current fusion predicts
$Z=\text{edema}\cup\text{scar}$ and obtains pure edema by subtracting scar at
fusion time. Further work will focus on strengthening edema supervision,
calibrating edema-zone fusion and adding inter-slice context.

% ===========================================================================
\section{Conclusion}
% ===========================================================================

We presented a coarse to fine heterogeneity-aware framework for myocardial
pathology segmentation with variable
sequence availability, heterogeneous labels and contrast-free cine scar evidence.
Its coarse to fine cascade first anchors the problem with cardiac anatomy and
then decodes pathology inside the anatomical prior. Presence-gated encoders with
a supervision-masked loss allow pooled training across heterogeneous sequence and
label availability, the cine branch uses temporal motion evidence, and the edema
expert predicts the injury-zone label used for fusion. On official validation, the best MoSAIC Dice entries were $0.6965$ for MyoPS
scar, $0.5983$ for MyoPS edema, and $0.2069$ for CineMyoPS scar.

\bibliographystyle{splncs04}
\bibliography{refs}

% ===========================================================================
\appendix
% ===========================================================================

\section{Thin-Plate Spline Alignment Field}
\label{app:tps}

With target control points $\{p_k\}$ on a regular $g\times g$ lattice spanning
$[-r,r]^2$ and predicted offsets $\theta^{(m)}_k$, the thin-plate
spline~\cite{bookstein1989principal} is the interpolant of
$\{p_k\mapsto\theta^{(m)}_k\}$ with minimal bending energy:
\begin{equation}
\mathcal{T}_{\theta}(q) \;=\; a_0 + A\,q \;+\; \sum_{k=1}^{K} w_k\,U\!\left(\lVert q-p_k\rVert\right),
\qquad U(d)=d^2\log d ,
\end{equation}
where $(w,a_0,A)$ solve the standard TPS linear system. That system depends only
on $\{p_k\}$, which are fixed, so it is inverted once and cached. The final
linear layer of the localization head $h^{(m)}$ is zero-initialized with its bias
set to the target lattice, giving a stable starting warp that is updated by the
task loss.

\section{Cross-Modality Gated Fusion}
\label{app:fusion}

The Stage-2a decoder fuses the aligned pyramids at level $\ell$ as
\begin{align}
a_\ell &= \tfrac{1}{2}\left[\sigma\!\left(\hat{f}^{(\Cz)}_\ell\right) + \sigma\!\left(\hat{f}^{(\TT)}_\ell\right)\right],\\
d_\ell &= \Phi_\ell\!\left(\left[\,\mathrm{up}(d_{\ell+1}),\;\; f^{(\LGE)}_\ell\odot a_\ell,\;\; \sigma\!\left(\hat{f}^{(\Cz)}_\ell\right),\;\; \sigma\!\left(\hat{f}^{(\TT)}_\ell\right)\right]\right),
\end{align}
initialized from the concatenated bottlenecks
$d_4=[f^{(\LGE)}_4,\hat{f}^{(\Cz)}_4,\hat{f}^{(\TT)}_4]$. When an auxiliary
sequence is absent, its masked feature pyramid is zero, so its sigmoid gate
channel is fixed at $\sigma(0)=\tfrac12$. This is a consequence of the gate
parameterization, not an identity gate; the absent sequence contributes no
image-derived feature map. The Stage-1 prior gates every decoder level,
$d_\ell \leftarrow d_\ell\odot\sigma(\mathrm{pool}(P))$.

\section{Auxiliary Training Objectives}
\label{app:losses}

\noindent\textbf{Cross-sequence myocardium consistency.} Three auxiliary heads
$g^{(m)}$ decode a myocardium mask from each sequence's own features. For the
available sequences $A=\{m:m_m=1\}$, each head is supervised against the
myocardium label and every available pair in
$\mathcal{Q}=\{(m,m'):m<m',\,m,m'\in A\}$ is driven into agreement:
\begin{equation}
\mathcal{L}_{\text{cons}}=\frac{\mathcal{L}_{\text{sup}}+
\mathcal{L}_{\text{pair}}}{N},\qquad N=|A|+|\mathcal{Q}|.
\end{equation}
Here $\mathcal{L}_{\text{sup}}$ sums the sequence-specific myocardium Dice
terms, and $\mathcal{L}_{\text{pair}}$ sums the pairwise Dice agreement terms.

\noindent\textbf{Stage-2b edema expert.} The reference pyramid is formed as
\begin{equation}
e^{\text{ref}}_\ell=\max\!\left(E^{\LGE}_\ell([\Img_\LGE,M]),\;E^{\Cz}_\ell([\Img_\Cz,M])\right),
\qquad
e^{\TT}_\ell=E^{\TT}_\ell([\Img_\TT,M]),
\end{equation}
with $M$ a dilated cardiac mask from Stage~1.

\end{document}
